\section{Measurement protocol}
\label{sec:protocol}

The claims we examine are quantitative, and several of them differ from what
we measure by amounts smaller than the run-to-run spread of a careless
experiment. What follows is therefore not boilerplate: each of the three
commitments below closes a specific way in which a number can be wrong without
anyone noticing.

\subsection{A pinned environment that refuses to be substituted}

Every measurement in this paper comes from one toolchain: Python
\PyVersion{}, PyTorch \TorchVersion{}, CUDA \CudaVersion{}, on an
\GpuName{} with driver \DriverVersion{}.

The pin is enforced rather than documented. Every Python process in the
repository is launched through a wrapper that inspects the interpreter and
its versions and exits non-zero when they do not match, so a measurement
cannot silently fall back to a different interpreter. This matters more for
state space models than for transformers: the selective scan runs an order of
magnitude slower without its compiled CUDA kernel, and a latency number
produced by that fallback looks like a result rather than a mistake.

The four experiments record their environment independently at run time, and
a test asserts that the four records agree on interpreter, framework, CUDA,
device, and driver. The sentence "one pinned environment" in the previous
paragraph is therefore checked rather than asserted; if any experiment were
re-run elsewhere, the build of this paper would fail instead of quietly
printing a version string that no longer describes the data.

\subsection{Numbers that cannot be typed by hand}

Every measurement in this paper is generated from the comma-separated files
committed alongside the code. Tables are emitted by a script that reads those
files; quantities that appear in running prose are emitted by the same script
as macros and referenced by name. The factorial effects in
Section~\ref{sec:factorial} are computed by calling the repository's own
analysis module, not by a second implementation written for the manuscript.

A test then reads every section source and fails if it finds a number that
looks like a measurement---two or more decimal places, or a percentage---that
is not one of a small set of explicitly justified exceptions. The exceptions
are values quoted from other papers, which are not ours to regenerate and are
attributed where they appear.

This is stricter than it may seem necessary. Transcription is exactly the
kind of error that survives proofreading, because a wrong digit is still a
plausible digit. The rule also removes a subtler failure: when a measurement
is repeated and the data changes, hand-written prose keeps the old value, and
nothing in the document indicates that it has gone stale.

\subsection{Predictions fixed before measurement}

Each experiment registered its predictions in a design document committed
before the corresponding measurement ran, and those predictions were not
edited afterwards.

Two of them did not survive contact with the data. The prediction about how
sharply object coverage degrades with object size turned out to depend on
the scale used to measure the degradation, and we describe in
Section~\ref{sec:dilution} why we moved the verdict to a scale-independent
statement rather than declaring the prediction met. The sanity band we
registered for the four training cells was exceeded from above by both
hierarchical cells, and Section~\ref{sec:factorial} reports what we did to
establish that the cause was not a defect in the pipeline.

We report both because this paper's thesis is that quantitative claims in
this area circulate without independent measurement. An argument of that
shape is worth very little from authors who filter their own results, and a
reader has no way to distinguish a study that registered three predictions
and reports three from one that registered five and reports the three that
worked.

\subsection{What the protocol does not buy}

A pinned environment makes a number reproducible; it does not make it
general. Every measurement here comes from a single consumer GPU with 8~GiB
of memory, and Section~\ref{sec:complexity} shows one claim whose failure to
reproduce cannot be separated from that constraint. The training experiment
uses one small dataset at low resolution. And the three architectures do not
expose a structurally identical point at which to probe internal behaviour,
which bounds how finely the receptive-field and coverage comparisons in
Sections~\ref{sec:erf} and~\ref{sec:dilution} can be read; we state that
bound where those results are used rather than only in the limitations.
